\section{Introduction}
\label{sec:introduction}

Multimodal judges compare vision--language outputs and provide preference
signals for model development. They are commonly summarized by static
preference accuracy: given an image, question, and two responses, how often is
the correct response selected? This cross-sectional score does not reveal how
the same decision changes when its visual evidence changes.

This question is adjacent to, but distinct from, controlled perturbation and
failure-to-update studies \citep{lee2026mmjudgebias,park2026mitigating,zhu2026visualflip}.
RewardLens evaluates the residual behavioral variation among judges that an
\emph{independent} static preference evaluation assigns equal or nearly equal
scores. It also tests both a relevant edit, for which the gold preference must
change, and an irrelevant edit, for which it must remain fixed.

We introduce \textbf{RewardLens}, a behavioral audit asking:
\emph{among judges with equal or near-equal measured static accuracy, how much
intervention-behavior variation remains?} An independent natural-image set
measures factor-specific static accuracy. A separate controlled audit then
holds the question and responses fixed while changing the rendered scene. Each
triplet contains a base image, a relevant edit that reverses the gold
preference, and an irrelevant edit that preserves it. Relevant Adaptation (RA)
measures correctness after the relevant edit, conditional on a correct base
decision; Irrelevant Invariance (II) measures the analogous correctness after
the irrelevant edit.

Across eight multimodal judges and four factors---Count, Attribute, Presence,
and Spatial---static similarity can coexist with different measured
intervention behavior. Phi-3.5-Vision and LLaVA-OneVision-7B tie at $80.5\%$
static Attribute accuracy, yet differ in RA by $15.1$ percentage points
(paired-bootstrap 95\% CI: $10.2$--$20.4$ pp). Across all seven within-factor
pairs matched within one static-accuracy point, the descriptive median RA gap
is $15.1$ points, and four gaps are at least $10$ points. Qwen3-VL-4B provides
a concrete transition-level example: on $166/200$ base-correct Count triplets,
the gold preference changes but its prediction does not.

RewardLens does not identify internal reasoning or replace static evaluation.
It shows that a shared measured static score can leave consequential behavioral
variation unresolved. Our contributions are a static-score-matched measurement
framing, a programmatically constructed and validated RewardLens-CLEVR audit
with relevant and irrelevant edits, and an eight-judge evaluation with paired
uncertainty and pre-specified matching-band checks.
